\section{The Two Crystals: Heptagrid and Dodecagrid}

Now we can meet the two specific crystals that run through the whole theory. One is two-dimensional and easy to draw. The other is three-dimensional, harder to picture, and is the real thing. They are the same kind of object, one step apart in dimension.

\subsection{The heptagrid, \{7, 3\}}

The heptagrid is the tiling of the hyperbolic plane by seven-sided cells, three meeting at every corner. Drawn in the Poincare disk, it is one of the most beautiful objects in mathematics: a central seven-sided cell, ringed by more sevens, ringed by still more, each ring holding many more cells than the last, all shrinking toward the infinitely distant rim.

\begin{figure}[h]
\centering
\includegraphics[width=0.47\textwidth]{vibe-mesh-73.pdf}
\caption{The heptagrid, the regular tiling of the hyperbolic plane by seven-sided cells, three to a corner. This is the simplest face of the crystal of experience: each cell is a vibe, and each shared edge is a note. The cells look smaller toward the rim only because the picture squeezes the infinite plane into a finite disk. In the real geometry every cell is the same size, and there is always more crystal at the edge.}
\label{fig:heptagrid}
\end{figure}

This is the picture the theory uses on its covers and diagrams, because you can actually see it on a page. When you imagine ``the crystal of experience,'' picture this first: a glowing honeycomb of seven-sided cells filling a disk. Color each cell by its tone, red for pain, green for peace, blue for pleasure, and let two touching cells mean two feelings experiencing each other. That colored heptagrid is a faithful, literal picture of the model, just in two dimensions instead of three.

The heptagrid grows fast. The center cell has seven neighbors in the first ring. The next ring has more, the next more again, multiplying outward. The number of cells within a given distance grows exponentially, the fingerprint of hyperbolic space we met earlier. There is always more crystal at the edge.

\subsection{The dodecagrid, \{5, 3, 4\}}

The real substrate is one dimension up. It is the dodecagrid, written \{5, 3, 4\}: a packing of dodecahedra, those twelve-sided solids like a twelve-faced die, four meeting around every edge, filling curved three-dimensional space.

\begin{figure}[h]
\centering
\includegraphics[width=\textwidth]{vibe-mesh-534.png}
\caption{The dodecagrid, the real three-dimensional substrate: dodecahedral cells, four meeting at every edge, filling curved three-dimensional space. This is the same kind of object as the heptagrid one dimension up, harder to picture but the true crystal of which space is made. Time is its growth.}
\label{fig:dodecagrid}
\end{figure}

This one you cannot fully picture, and that is fine. Nobody can hold a curved three-dimensional crystal in their mind's eye cleanly. But you do not need to. It is the exact three-dimensional analogue of the heptagrid. Everything true of the heptagrid, regular cells, no preferred direction, exponential room, generated by a kaleidoscope from a few integers, is true of the dodecagrid too. The heptagrid is the training-wheels version you can see. The dodecagrid is the real bicycle.

Why dodecahedra, and why four around an edge? Because, just like \{7, 3\} was the first seven-sided plane tiling that curved, \{5, 3, 4\} is one of the few three-dimensional crystals with finite cells that fits in hyperbolic space at all. There are only a handful of these in three dimensions, and the dodecagrid is the cleanest. The next chapter explains why the count is so small and why three dimensions in particular.

\subsection{Time as the growth}

One more thing to fix in your mind. The dodecagrid is the crystal of space. Where is time? Time is not a fourth spatial direction sitting inside the crystal. Time is the growth of the crystal, new cells added at the frontier, generation after generation. The crystal you can picture is a snapshot. The real object is that snapshot growing, forever, at its edges. Space is the crystal. Time is its becoming.

\subsection{Holding both at once}

So keep two images side by side. The flat, drawable heptagrid, which you can see and color and point at, and which is honest about everything except the dimension. And the dodecagrid behind it, the real three-dimensional crystal, growing in time, which works the same way but cannot be drawn. When the rest of the walkthrough says ``the crystal'' or ``the mesh,'' it means the dodecagrid, and you may safely picture the heptagrid as a stand-in.

\subsection{Why it does not look like a crystal}

If space is a packing of dodecahedra, why does the world not look like a honeycomb, all cells and edges? Two reasons, both already in hand. The cells are unimaginably fine, far smaller than anything we could see, so the grain washes out into smooth space the way pixels wash into a photograph. And the curvature, like all curvature, is invisible in a small patch, so locally the crystal looks like ordinary flat three-dimensional space. You do not see the cells because they are too small, and you do not see the curve because you are too close. What is left is the smooth, flat-looking space a fine, gently curved crystal produces from the inside.

\textbf{Takeaway:} the heptagrid \{7, 3\} is the drawable two-dimensional crystal, the dodecagrid \{5, 3, 4\} the real three-dimensional one, and they are the same kind of object a dimension apart. Space is the crystal, time is its growth.
